\documentclass[10pt,a4paper]{amsart}
\usepackage[margin=19mm]{geometry}
\usepackage{amsmath,amssymb,mathtools}
\usepackage[scr=rsfs,cal=euler]{mathalfa}
\usepackage{xfrac}
\usepackage[unicode,hidelinks,bookmarks=false]{hyperref}
\newcommand{\ie}{i.e.\ }
\newcommand{\cf}{cf.\ }
\newcommand{\resp}{respectively\ }
\newcommand{\Subsection}{\subsection}
\providecommand{\nobreakdash}{\nobreak\hbox{-}}
\setlength{\emergencystretch}{2em}
\multlinegap0pt
\usepackage{bm}
\usepackage{bbm}
\usepackage{dsfont}
\usepackage{xspace}
\usepackage{cancel}
\usepackage{mathtools}
\usepackage{amsthm}
\makeatletter
\@ifundefined{thm}{\newtheorem{thm}{Thm}}{}
\@ifundefined{lemma}{\newtheorem{lemma}[thm]{Lemma}}{}
\makeatother
\newcommand{\Invt}{\widetilde{\mathrm{Inv}}}
\newcommand{\Inv}{{\mathrm{Inv}}}
\newcommand{\inc}{\subseteq}
\title[Generalised flip order on the faces of nestohedra]{Generalised flip order on the faces of nestohedra}
\author{Pierre-Louis Curien, B\'er\'enice Delcroix-Oger, Jovana Obradovi\'c}
\date{Source: arXiv:2607.20132v2 (CC BY 4.0)}
\begin{document}
\maketitle

\noindent\emph{Source and licence.} Generalised flip order on the faces of nestohedra. Version arXiv:2607.20132v2 (CC BY 4.0), distributed under the Creative Commons Attribution 4.0 International licence, as recorded in the arXiv licence metadata for this exact version and shown on the abstract page https://arxiv.org/abs/2607.20132v2. Full rights notice: licenses/arxiv-2607.20132-CC-BY-4.0.txt.\par\smallskip

\noindent\textit{Editorial scope.} Counted unit: the Lemma whose source label is \texttt{s14} in the article (source0.tex lines 2157--2159, source label \texttt{s14}), delivered together with the article's own complete original proof of it (source0.tex lines 2161--2253, one \texttt{proof} environment of 93 lines). No mathematics is written, repaired or re-derived: statement and proof are the article's own text line for line. Required context retained uncounted: s000 (lines 2106--2108); s00 (lines 2113--2115). Every definition, display and label of those lines is retained unchanged. Omitted from this package and not redistributed: 1--2105, 2109--2112, 2116--2156, 2160--2160, 2254--4739. Nothing inside a retained excerpt is omitted or altered: every retained body is delivered exactly as the article's lines are written, so no omission receipt of any kind accompanies this package. Citations: the retained excerpts carry no citation command at all, so no bibliography entry of the article is transcribed and no reference is redistributed. Reference adaptation: none was needed and none was made; every reference inside the delivered unit resolves inside it, so no unresolved reference or citation is produced. Printed slips kept literal: no printed word, symbol, hypothesis or number of the counted statement or of its proof was altered. Layout and class: the article's own class is \texttt{article} with packages including amssymb, amsthm, amsmath, bm, xcolor, tikz, hyperref, todonotes, stmaryrd, pifont, float, caption, centernot, ulem; the wrapper is the lane's pinned \texttt{amsart} editorial wrapper at 10pt on A4, which restates the theorem-like environments the retained text needs and adds the article's own macro definitions that the retained text needs (\texttt{\textbackslash Inv}, \texttt{\textbackslash Invt}, \texttt{\textbackslash inc}), with extra wrapper packages (bm, bbm, dsfont, xspace, cancel, mathtools). The delivered numbering is the unit's own. Assets: no retained span contains a figure environment, a graphics-inclusion command, a PDF-page inclusion, a table or a TikZ picture; no image file is copied into this package, the package declares no asset dependency and no image file is redistributed with it. Licence: version arXiv:2607.20132v2 of this article is distributed under the Creative Commons Attribution 4.0 International licence, as recorded in the arXiv licence metadata for this exact version and shown on the abstract page https://arxiv.org/abs/2607.20132v2. A full rights notice is delivered at licenses/arxiv-2607.20132-CC-BY-4.0.txt. Characters: every retained excerpt is pure ASCII, so the delivered engine, XeLaTeX with its default Latin Modern text font, reports no missing character.
\begin{lemma}\label{s000}
If ${\bf H}$ is an ordered right-filled hypergraph and if $K\subseteq H$, then the restricted hypergraph ${\bf H}_{K}$ is also right-filled (with respect to the induced ordering on $K$).
\end{lemma}

\begin{lemma}\label{s00}
Suppose that ${\bf H}$ is an ordered right-filled hypergraph, that let $K\inc H$ is connected in ${\bf H}$ and that, for some $a,b,c\in H$, it holds that $a,c\in K$, $b\not\in K$ and $a<b<c$. Then $K\cup \{b\}$ remains connected in ${\bf H}$.
\end{lemma}

\begin{lemma}\label{s14}
Let ${\bf H}$ be an ordered right-filled hypergraph and let $S,T:{\bf H}$ be such that $\Inv_{\bf H}(S)= \Inv_{\bf H}(T)$ and $\Invt_{\bf H}(S)\subsetneq \Invt_{\bf H}(T)$. Then there exists a construct $U:{\bf H}$ such that $U\lessdot^+ T$, $\Inv_{\bf H}(S)= \Inv_{\bf H}(U)$ and $\Invt_{\bf H}(S)\subseteq\Invt_{\bf H}(U)\subsetneq \Invt_{\bf H}(T)$. 
\end{lemma}

\begin{proof}
 Pick $(a,b)\in \Invt_{\bf H}(T)\backslash \Invt_{\bf H}(S)$ and let $X=\{x_1<\cdots <x_n\}$ be the vertex of $T$ such that $a=x_i$ and $b=x_j$ for some $1\leq i<j\leq n$. 

 

We define $U:=T[B\langle A\rangle /X]$, 
 where the choice of decomposition $X=B\cup A$ is such that $\{x_1,\dots,a\}\subseteq A$, $\{b,\cdots,x_n\}\subseteq B$ and for $x\in \{x_{i+1},\dots,x_{j-1}\}$, we set $x\in A$ if $(a,x)$ is a neutral pair of $S$; otherwise, we set $x\in B$. By construction, $(a,b)\not\in \Invt_{\bf H}(U)$. In addition, no new neutral pairs are produced by the squashing that defines $U$, so we have that $\Invt_{\bf H}(U)\subsetneq \Invt_{\bf H}(T)$. 

\smallskip

Suppose that the subconstruct of $T$ rooted at $X$ is given by $X(C_1,\dots,C_t):{\bf K}$, where ${\bf K},X\leadsto K_1,\cdots ,K_t$ and $C_i:{\bf K}_i$ and let $A^1<\cdots <A^q$ be the disintegration of $A$ involved in the $(B,A)$-squashing of $X$ that defines $U$. For each $1<j<q$, write $I_j\subseteq \{1,\dots,t\}$ for the set of indices such that, for $i\in I_j$, the set $A^j\cup K_i$ is not connected in ${\bf K}\backslash B$ and $A^j<K_i$. In order to show that $U$ has other desired properties, we use the following three claims.
\begin{itemize}
\item[($\heartsuit$)] For each $1\leq k<l\leq n$, if $(x_k,x_l)$ is neutral in $S$, then, for all $k\leq k' < l' \leq l$, $(x_{k'},x_{l'})$ is also neutral in $S$.
\item[($\clubsuit$)] If $(x,y)\in A^r\times A^s$ for some $1<r<s<q$, then $(x,y)$ is not neutral in $S$.
\item[($\spadesuit$)] If $(u,v)\in A^j\times K_i$ for some $1<j<q$ and $i\in I_j$, then $(u,v)$ is not good in $S$.
\end{itemize}
 
We first note that the property ($\heartsuit$) implies that $A<B$. Indeed, if there would exist $k<l$ in the interval $[i+1,j-1]$ such that $x_k\in B$ and $x_l\in A$, then, by the definition of $U$, we would have that $(a,x_l)$ is neutral in $S$ while $(a,x_k)$ is not, but this is impossible by ($\heartsuit$), as it requires $(a,x_k)$ to also be neutral in $S$. Having established that $A<B$, we conclude that $T$ can be recovered from $U$ by successive fusions of the $(A^i)$'s with their parent node, which entails that $U\lessdot^{+} T$.

 Next, the properties ($\heartsuit$) and ($\clubsuit$) guarantee that no neutral pair of $S$ gets abolished in $U$, i.e., that $\Invt_{\bf H}(S)\subseteq\Invt_{\bf H}(U)$. To see this, we observe that, having that $\Invt_{\bf H}(S)\subseteq\Invt_{\bf H}(T)$, there are two scenarios in which a neutral pair $(x,y)$ of $S$ would hypothetically be abolished in $U$:
\begin{itemize}
\item $x\in A_r$ and $y\in A_s$ for some $1<r<s<q$: this is impossible by ($\clubsuit$);
\item $x\in A$ and $y\in B$ - in this case, we further distinguish subcases relative to the exact position of $x$ and $y$ in $X$:
\begin{itemize}
\item if $x\leq a$ and $y\geq b$, then, by ($\heartsuit$), $(a,b)$ would have to be neutral in $S$, which contradicts the assumption that $(a,b)\in \Invt_{\bf H}(T)\backslash \Invt_{\bf H}(S)$;
\item if $x\leq a$ and $y=x_k$ for some $k\in[i+1,j-1]$, then the contradiction arises as the definition of $B$ implies that $(a,y)$ is not neutral in $S$, while, at the same time, ($\heartsuit$) requires $(a,y)$ to be neutral in $S$ (since $(x,y)$ is neutral in $S$);
\item if $x=x_k$ for some $k\in[i+1,j-1]$ and $y\geq b$, then the definition of $A$ implies that $(a,x)$ is neutral in $S$ and ($\heartsuit$) implies that $(x,b)$ is also neutral in $S$ (since $(x,y)$ is neutral in $S$), which, altogether, by transitivity, implies that $(a,b)$ is neutral in $S$, contradicting once again the assumption that $(a,b)\in \Invt_{\bf H}(T)\backslash \Invt_{\bf H}(S)$;
\item if $x=x_k$ and $y=x_l$ for some $k<l$ in the interval $[i+1,j-1]$, then, by the definitions of $A$ and $B$, we have that $(a,x)$ is neutral in $S$ and $(a,y)$ is not neutral in $S$, while, at the same time, the transitivity of neutrality implies that $(a,y)$ must be neutral in $S$ (since $(x,y)$ is neutral in $S$), contradiction. 
\end{itemize} 
\end{itemize}

 Finally, the property ($\spadesuit$) guarantees that no good pair of $S$ gets abolished in $U$, i.e., that $\Inv_{\bf H}(S)=\Inv_{\bf H}(U)$. 

\smallskip

 The lemma therefore follows by proving the properties ($\heartsuit$), ($\clubsuit$) and ($\spadesuit$). 
\smallskip

{\em The proof of} ($\heartsuit$). Suppose that there exist indices $1\leq k\leq k'<l'\leq l\leq n$ such that $(x_k,x_l)$ is neutral in $S$ and $(x_{k'},x_{l'})$ is not neutral in $S$. Let $V$, $Y$ and $Z$ be the vertices of $S$ such that $x_k,x_l\in V$, $x_{k'}\in Y$ and $x_{l'}\in Z$, and let $W$ be the first common ancestor of $Y$ and $Z$ in $S$. Then the condition $\Inv_{\bf H}(S)= \Inv_{\bf H}(T)$ implies that the subtrees of $S$ rooted at $V$ and $W$ are disjoint, but this is in contradiction with Lemma \ref{s00}.

\smallskip 
{\em The proof of} ($\clubsuit$) {\em and} ($\spadesuit$). Having defined $$N:=\{(x,y)\in A_i\times A_j\,|\, i<j\}\enspace \mbox{ and } \enspace M:=\{(x,y)\in A^j\times K_i\,|\, 1<j<q, i\in I_j\},$$ our goal is to show that $\Invt_{\bf H}(S)\cap N=\Inv_{\bf H}(S)\cap M=\emptyset$. Suppose, for contradiction, that there exists $(x,y)\in (\Invt_{\bf H}(S)\cap N)\cup (\Inv_{\bf H}(S)\cap M)$. Let $Y$ be the node of $S$ such that $x\in Y$, and suppose that the subconstruct of $S$ rooted at $Y$ is given by $Y(T_1,\dots,T_r):{\bf G}$, where ${\bf G},Y\leadsto G_1<\cdots <G_r$ and $T_i:{\bf G}_i$. Since ${\bf G}$ is connected, there exists a chain of hyperedges $E_1,\dots,E_m\in{\bf G}$, such that $x\in E_1$, $y\in E_m$, and there exists a choice of $q_i\in E_i\cap E_{i+1}$, $1\leq i\leq m-1$. Our proof proceeds by induction on $m$.

\smallskip

Suppose that $m=1$. Since $x$ and $y$ are separated in $U$, i.e., belong to different connected components of ${\bf K}\backslash B$, $E_1$ must intersect at least one vertex on the path, in $U$, starting from $B$ and going to the root vertex of $U$. We proceed by analysing the cases relative to whether $(x,y)$ is a neutral or a good pair of $S$. We shall work out the details only for the case that it is neutral, with the indication that the proof strategy is the same in the case of a good pair. 
\begin{itemize}
\item Suppose that there exists an $e\in E_1\cap B$. Since no element of $B$ can appear strictly above $Y$ in $S$, as this would produce a good pair not present in $T$, it must be the case that $e\in Y$. We now proceed by analysing the position of $x,y$ and $e$ in $X$, relative to $a$ and $b$.
\begin{itemize}
\item if $e<b$, then, by definition of $B$, we have that $(a,e)$ is not neutral in $S$, and hence $a$ must appear either above $Y$, or below $Y$ or on a branch disjoint from $Y$; the first two possibilities are in contradiction with the fact that $\Inv_{\bf H}(S)= \Inv_{\bf H}(T)$, and the third one with the right-filled property in the form of Lemma \ref{s00};
\item if $e\geq b$ and $x<y\leq a$ or $x\leq a<y$, then, by ($\heartsuit$), since $(x,e)$ is neutral in $S$, $(a,b)$ would also have to be neutral in $S$, which contradicts the starting assumption;
\item if $e\geq b$ and $a\leq x<y$, then, by definition of $A$, we have that $a\in Y$, and hence, it must be the case that $e>b$ with $b$ appearing either below $Y$ or on a branch disjoint from $Y$; the first possibility is in contradiction with the fact that $\Inv_{\bf H}(S)= \Inv_{\bf H}(T)$ and the second one with the right-filled property in the form of Lemma \ref{s00}; 
\end{itemize}
\item If $E_1\cap B=\emptyset$, let $V$ be any vertex of $U$ below $B$ such that $E_1\cap V\neq \emptyset$. Observe that $V$ is then also a vertex of $T$ below $X$, and moreover that the path from $V$ to the parent of $B$ in $U$ is the same as the path from $V$ to the parent of $X$ in $T$. Pick an element $e\in E_1\cap V$. Note that if $e<y$, then $(e,y)$ is a good pair in $T$ and hence in $S$, and hence $e$ must appear strictly below $Y$ is $S$, contradicting the fact that $e\in E_1\subseteq G$. Likewise, if $e>y$, then $(e,y)$ is not a good pair in $T$ and hence can neither be good in $S$. Therefore, since $e\in E_1\subseteq G$, the only possibility is that $e\in Y$, but this is contradicting the fact that $\Invt_{\bf H}(S)\subsetneq \Invt_{\bf H}(T)$. 
\end{itemize}
This finishes the proof for the case $m=1$.

\smallskip

Suppose now that $m>1$ and consider $q_{m-1}$. We differentiate two cases.
 
\smallskip
\indent Case (1). If $q_{m-1}<x$, then, by applying the right-filled property of ${\bf G}$ (ensured by Lemma \ref{s000}) to $E_{m}$, we conclude that there exists a hyperedge $E$ of ${\bf G}$ such that $x,y\in E$, which takes us back to our base case.


 
\smallskip

 
\indent Case (2). If $q_{m-1}>x$, then $(x,q_{m-1})$ is either a neutral or a good pair of $S$. Suppose that $(x,y)$ is neutral in $S$.
\begin{itemize}
\item If $q_{m-1}$ is strictly above $Y$, and
\begin{itemize}
\item if $x<y<q_{m-1}$, then $(x,q_{m-1}),(y,q_{m-1})\in \Inv_{\bf H}(S)$, and hence, since ($\spadesuit$) holds by induction, $(x,q_{m-1}),(y,q_{m-1})\in \Inv_{\bf H}(U)$, but this configuration is impossible as it produces a cycle in the tree $U$ (given that $x$ and $y$ are separated there);
\item if $x<q_{m-1}<y$, then $(x,q_{m-1})\in \Inv_{\bf H}(S)$, and hence, by induction, $(x,q_{m-1})\in \Inv_{\bf H}(U)$, which means that $q_{m-1}$ and $y$ are separated in $U$, which is impossible by our base case (applied to the pair $(q_{m-1},y)$ and the hyperedge $E_m$ containing them).
\end{itemize}


\item Suppose that $q_{m-1}\in Y$, and hence also $q_{m-1}\in X$. We first show that, in fact, $q_{m-1}\in A$. Suppose, for contradiction, that $q_{m-1}\in B$. Then, if $q_{m-1}\geq b$ and
\begin{itemize}
\item if $x<y<a$ or $x<a<y$, we get a contradiction with the fact that $(a,b)$ is not neutral in $S$, using ($\heartsuit$);
\item if $a<x<y$, by definition of $A$, we get that $a\in Y$, and hence that $(a,q_{m-1})$ is neutral in $S$, which, using ($\heartsuit$) once again, leads to a contradiction with the fact that $(a,b)$ is not neutral in $S$;
\end{itemize}
while, if $q_{m-1}< b$, by definition of $B$, we have that $(a,q_{m-1})$ is not neutral in $S$ and hence
\begin{itemize}
\item if $x<y<a$ or $x<a<y$, we get a contradiction once again using ($\heartsuit$);
\item if $a<x<y$, then, by definition of $A$, we get that $a\in Y$, i.e., that $(a,x)$ is neutral, but since $(x,q_{m-1})$ is also neutral, we get a contradiction by transitivity of neutrality. 
\end{itemize}

Now, since $q_{m-1},y\in Y$ and $q_{m-1},y\in E_m$, our base case implies that $q_{m-1}$ and $y$ cannot be separated in $U$, i.e., that $q_{m-1}\in A_j$. Therefore, $x$ and $q_{m-1}$ are separated in $U$, and this is impossible by induction.
\end{itemize}
This concludes the proof for the case $(x,y)\in \Invt_{\bf H}(S)$. We proceed analogously if $(x,y)\in \Inv_{\bf H}(S)$, in which case the induction argument uses ($\clubsuit$).
 \end{proof}

\end{document}
